% Generated draft: requires mathematical and source review.
\documentclass[11pt]{article}
\usepackage[margin=25mm]{geometry}
\usepackage{fontspec}
\setmainfont{Latin Modern Roman}
\usepackage{amsmath,amssymb,mathtools,unicode-math}
\setmathfont{Latin Modern Math}
\usepackage{enumitem,needspace,xurl,hyperref}
\setlength{\parindent}{0pt}
\setlength{\parskip}{5pt}
\setlength{\emergencystretch}{3em}
\newcommand{\N}{\mathbb{N}}
\newcommand{\Z}{\mathbb{Z}}
\newcommand{\Q}{\mathbb{Q}}
\newcommand{\R}{\mathbb{R}}
\newcommand{\C}{\mathbb{C}}
\newcommand{\F}{\mathbb{F}}
\newcommand{\A}{\mathbb{A}}
\newcommand{\PP}{\mathbb{P}}
\newcommand{\E}{\mathbb{E}}
\newcommand{\eps}{\varepsilon}
\newcommand{\dd}{\,\mathrm{d}}
\begin{document}
\section*{UnsolvedMath ID: 20003347}
\textbf{Source problem code:} catalogue entry
\subsection{Definitions and mathematical statement}
The source records a historical exception to an unreproduced Conjecture 48(3). Its displayed tuple denotes the subgroup
\[
 D=\{\operatorname{diag}(s,s^{-1},t,t^{-1}):s,t>0\}
 =\{\operatorname{diag}(e^u,e^{-u},e^v,e^{-v}):u,v\in\R\}
\]
in \(\operatorname{SL}_4(\R)\). For \(X_4=\operatorname{SL}_4(\R)/\operatorname{SL}_4(\Z)\), a divergent \(D\)-orbit means that its orbit map is proper, equivalently escaping every compact subset as its parameter leaves every compact subset of \(D\). The technical adjective ``obvious'' is not defined in this excerpt and cannot be replaced by an invented criterion.

The aid attributes to Tamam's published Example 4.1 and Theorem 1.3 the existence of non-obvious divergent \(D\)-orbits, resolving the historical exceptional subgroup. The original notes say that Weiss's 2004 result [200] covered the other \(\operatorname{SL}_4\) diagonal subgroups. This is not a reconstruction of every assertion in the missing Conjecture 48.

The supplied local calculation can be stated independently. Restriction of the \(\operatorname{SL}_4\) roots to logarithmic coordinates \((u,v)\) gives the \(C_2\)-shaped set
\[
 \{\pm2u,\pm2v,\pm(u+v),\pm(u-v)\},
\]
with the cross weights occurring twice and \(\pm2u,\pm2v\) once each. Let \(E=\operatorname{span}(e_1,e_2)\), \(F=\operatorname{span}(e_3,e_4)\). For a real two-plane \(W\), choose its nonzero decomposable Pl\"ucker vector
\[
 p_W=\sum_{i<j}p_{ij}e_i\wedge e_j\in\Lambda^2\R^4,
 \qquad p_{12}p_{34}-p_{13}p_{24}+p_{14}p_{23}=0.
\]
The zero-weight space is \(\operatorname{span}(e_1\wedge e_2,e_3\wedge e_4)\). The aid's exact classification is
\[
 \operatorname{proj}_{(\Lambda^2\R^4)^D}(p_W)=0
 \Longleftrightarrow p_{12}=p_{34}=0
 \Longleftrightarrow W=L_E\oplus L_F
\]
for lines \(L_E\subset E\), \(L_F\subset F\).

Among those split planes, there is a nontrivial ray
\[
 d_{r,s}(\tau)=\operatorname{diag}(e^{r\tau},e^{-r\tau},e^{s\tau},e^{-s\tau}),
 \qquad\tau\geq0,\quad(r,s)\ne(0,0),
\]
with \(d_{r,s}(\tau)p_W\to0\) exactly when at least one of the two factor lines is a coordinate eigenline: \(L_E=\R e_1\) or \(\R e_2\), or \(L_F=\R e_3\) or \(\R e_4\). The remaining wedge weights are \(u+v,u-v,-u+v,-u-v\). This statement concerns real decomposable vectors in one fundamental representation; it does not classify arithmetic lattice trajectories or every mechanism called obvious divergence.

The later source paragraph is a separate cone discussion: it reports cones with non-obvious divergent trajectories, including Weyl chambers, and a rank-two \(\operatorname{SL}_3\) example with only obvious divergence whose argument used dimension two [200]. The aid reports a subsequent rank-two cone classification by Tamam and a 2024 Datta--Tamam preprint showing obviousness for every divergent orbit of a specified fundamental-weight cone \(S^+\). The definition of that cone and all theorem hypotheses are absent, so no assertion for all cones is made. These later literature assessments, the subgroup exception, and the local Pl\"ucker theorem remain distinct.
\subsection{Short English statement}
The supplied literature assessment resolves the historical exceptional diagonal subgroup. Separately, the local calculation identifies which decomposable two-plane vectors have zero fixed projection and which can contract along a ray. Neither statement supplies a full definition or classification of every cone-divergence notion in the missing conjecture.
\subsection{Sources}
\begin{itemize}
\item Local UnsolvedMath registry, permanent ID \(20003347\), exact source code \texttt{catalogue entry}; canonical \texttt{aim-other-notes.json}, \texttt{notes[153]}, AIM domain Other, registry category Miscellaneous, tag \texttt{conjecture}.
\item AIM workshop ``Emerging applications of measure rigidity,'' historical Conjecture 48(3) and following cone discussion: \url{https://aimath.org/WWN/measrigid/measrigid.pdf}. The aid identifies [200] with Weiss's 2004 result. No PDF was fetched.
\item Tamam, published Example 4.1 and Theorem 1.3, and a rank-two cone classification; Datta--Tamam, 2024 preprint, for the specified \(S^+\) cone, as cited by the aid. No paper titles, identifiers, full authoritative entries, or definition of \(S^+\) are supplied. The original literature field is empty.
\item Research attempt 2: result \texttt{partial\_result}, problem status at run \texttt{solved}, rigor label \texttt{proved}, novelty confidence \texttt{low}; classification \texttt{SOLVED-IN-LITERATURE}, registry status \texttt{solved}. Created and updated 14 August 2026. These are supplied assessments, not a blanket resolution of unreproduced cone conjectures.
\item Attribution: Ulam AI / UnsolvedMath Contributors, CC BY 4.0. The machine-generated research aid is not peer reviewed and requires independent verification before citation.
\end{itemize}
\end{document}
